% =====================================================================
%  LA CRYPTE DE L'ORDRE DES NOMBRES PREMIERS
%  Sujet ludique — Mathématiques expertes, terminale générale
%  Quête de calcul matriciel : somme, différence, produit, inversion,
%  comparaison de matrices 2x2, 3x3, 4x4 — chiffrement modulo 7 et 13.
%  Compilation : pdflatex (2 passes) ou Overleaf.
% =====================================================================
\documentclass[11pt,a4paper]{article}

\usepackage[utf8]{inputenc}
\usepackage[T1]{fontenc}
\usepackage[french]{babel}
\usepackage{amsmath,amssymb}
\usepackage[margin=2.1cm,bottom=2.6cm]{geometry}
\usepackage{booktabs,array}
\usepackage{enumitem}
\usepackage{xcolor}
\usepackage{tcolorbox}
\usepackage{fancyhdr}


% ---------- Couleurs de la quête ----------
\definecolor{encre}{RGB}{26,32,44}
\definecolor{or}{RGB}{153,110,20}
\definecolor{argent}{RGB}{95,112,130}
\definecolor{emeraude}{RGB}{4,120,87}
\definecolor{jade}{RGB}{10,90,74}
\definecolor{parche}{RGB}{253,249,236}
\definecolor{nuit}{RGB}{24,28,44}
\definecolor{sang}{RGB}{150,32,32}

% ---------- Boîtes ----------
\newtcolorbox{missive}[1]{colback=parche,colframe=or,title={\textbf{#1}},
	coltitle=white,fonttitle=\normalsize,boxrule=1pt,arc=2mm,left=3mm,right=3mm}
\newtcolorbox{rituel}[1]{colback=white,colframe=argent,title={\textbf{#1}},
	coltitle=white,fonttitle=\normalsize,boxrule=1pt,arc=2mm,left=3mm,right=3mm}
\newtcolorbox{piege}{colback=sang!4,colframe=sang!80,boxrule=0.8pt,arc=1mm,
	left=3mm,right=3mm}
\newtcolorbox{conseil}{colback=emeraude!5,colframe=emeraude!70,boxrule=0.8pt,
	arc=1mm,left=3mm,right=3mm}
\newtcolorbox{precedemment}[1]{colback=argent!6,colframe=nuit!75,
	title={\textbf{#1}},coltitle=white,fonttitle=\normalsize,boxrule=1pt,
	arc=2mm,left=3mm,right=3mm}

% ---------- En-têtes ----------
\pagestyle{fancy}
\fancyhf{}
\fancyhead[L]{\small\textsc{La Crypte de l'Ordre des Nombres Premiers}}
\fancyhead[R]{\small\textit{Maths expertes, terminale générale}}
\fancyfoot[L]{\small\textit{Quête de calcul matriciel}}
\fancyfoot[C]{\small\thepage}
\fancyfoot[R]{\small\textit{Récompense de la quête : /20}}
\renewcommand{\headrulewidth}{0.6pt}

\setlength{\parindent}{0pt}
\setlist{itemsep=2pt,topsep=3pt}

% ---------- Raccourcis ----------
\newcommand{\Z}{\mathbb{Z}}
\newcommand{\Id}[1]{I_{#1}}
\newcommand{\Epx}[1]{\textcolor{or}{\textbf{[#1 XP]}}}

\begin{document}
	
	% =====================================================================
	\begin{center}
		{\LARGE\bfseries\textcolor{nuit}{LA CRYPTE DE L'ORDRE DES NOMBRES PREMIERS}}\\[2mm]
		{\large\itshape Une quête de calcul matriciel pour apprenti·e cryptographe}\\[2mm]
		\textbf{Mathématiques expertes, terminale générale}\quad\textbullet\quad
		Feuilleton en deux épisodes de 50 min\quad\textbullet\quad Barème : 20 XP (10 + 10)
	\end{center}
	
	\vspace{1mm}
	\begin{missive}{Message retrouvé dans le tombeau du dernier Grand Maître}
		\itshape
		Apprenti·e, si tu lis ces lignes, la Crypte t'a accepté·e.
		Des siècles durant, notre Ordre a protégé son trésor derrière trois sceaux
		de matrices. Les nombres premiers en sont l'alphabet, et les restes,
		seuls, disent la vérité : \textbf{seul le reste parle}.
		
		Trois épreuves t'attendent, de plus en plus profondes :
		le \textcolor{argent}{\textbf{Sceau d'argent}} (matrices $2\times 2$, rituel modulo 7),
		le \textcolor{emeraude}{\textbf{Sceau d'émeraude}} (matrices $3\times 3$, rituel modulo 13)
		et le \textcolor{jade}{\textbf{Sanctuaire de jade}} (matrices $4\times 4$, rituel modulo 13).
		
		Chaque épreuve livrera un fragment du message perdu de l'Ordre.
		Assemble-les, et le nom du trésor t'apparaîtra. Mais n'espère pas tout
		traverser en un jour : la Crypte ne s'ouvre que deux fois, cinquante
		minutes à chaque lune, et elle se referme entre les deux plongées.
		Le coffre final ne s'ouvrira qu'avec le code des premiers.
		Que la quête commence !
	\end{missive}
	
	% =====================================================================
	\section*{\textcolor{nuit}{Générique : règles de la quête}}
	
	{\small\itshape (Le générique passe avant chaque épisode : ces rituels
		valent pour toute la saison.)}
	
	\begin{rituel}{Les quatre rituels de l'Ordre}
		\begin{enumerate}[label=\textbf{R\arabic*.},leftmargin=10mm]
			\item \textbf{Seul le reste parle.} Dans la Crypte, tout est gravé en
			\emph{nombres premiers}. Chaque nombre premier $p$ gravé sur la pierre
			doit d'abord être réduit : on ne retient que son reste dans la division
			euclidienne par le module du sceau (7 ou 13).
			Par exemple, dans le sanctuaire (module 13) : $61 \equiv 9 \pmod{13}$,
			car $61 = 4\times 13 + 9$.
			\item \textbf{Les matrices vivent modulo $n$.} Deux matrices sont
			\emph{égales dans le sceau} si leurs coefficients sont égaux modulo $n$
			(7 ou 13). Tous les produits de matrices se réduisent modulo $n$.
			\item \textbf{L'alphabet du chiffre.} Les messages sont écrits en lettres
			majuscules non accentuées (l'Ordre a perdu ses accents : É se lit E).
			\begin{itemize}
				\item Dans les sceaux \textbf{modulo 13}, chaque nombre $x$
				($0\leqslant x \leqslant 12$) code une lettre :
				$\text{A}=0,\ \text{B}=1,\ \dots,\ \text{M}=12$
				(alphabet réduit de l'Ordre, annexe B).
				\item Dans le sceau \textbf{modulo 7}, chaque lettre est portée par un
				\emph{couple} $(a,b)$ de nombres ($0\leqslant a,b\leqslant 6$) :
				le rang de la lettre ($\text{A}=0$, $\dots$, $\text{Z}=25$) vaut
				$a + 7b$ (annexe C).
			\end{itemize}
			\item \textbf{La Table de l'Ordre.} Pour ouvrir le coffre final, chaque lettre
			vaut un nombre premier selon la Table de l'Ordre (annexe A) :
			$\text{A}=2,\ \text{B}=3,\ \text{C}=5$, etc.
		\end{enumerate}
	\end{rituel}
	
	\begin{conseil}
		\textbf{Mémo du cryptographe} (annexe D) : inverse d'une matrice $2\times 2$,
		produit de matrices, et le grand théorème des voleurs de clés :
		$(M N)^{-1} = N^{-1} M^{-1}$. Les annexes restent à ta disposition
		lors des deux plongées.
	\end{conseil}
	
	% =====================================================================
	\section*{\textcolor{nuit}{Épisode 2 : Le Nom de Pierre}}
	{\small\textbf{Deuxième séance : 50 min}\quad\textbullet\quad
		\textbf{Récompense : 10 XP}\quad\textbullet\quad
		Acte I : le réveil 5 min\quad\textbullet\quad Acte II : le mot de
		l'émeraude 10 min\quad\textbullet\quad Acte III : le Sanctuaire de jade
		25 min\quad\textbullet\quad générique de fin : le coffre 10 min}
	
	% ---------------------------------------------------------------------
	\begin{precedemment}{Précédemment, dans la Crypte... (résumé de l'épisode 1)}
		\itshape
		La première pierre est tombée. Dans la Porte d'argent, le héros a
		assemblé la pierre de voûte $K_1 = A + B - C$, démasqué le faux gravage,
		forgé sa première clé $K_1^{-1}$ grâce à l'identité de Bézout
		$3\times 5 - 2\times 7 = 1$, et arraché à la pierre le premier fragment
		du message : \textbf{« LA »}.
		
		Puis, devant le Sceau d'émeraude, il a calculé le sceau $K_2 = MN$,
		démasqué le parchemin du traître ($M'N'$ le faux, $N'M'$ le vrai) et mis
		la main sur la clé de décryptage $K_2^{-1}$. Mais la crypte s'est
		refermée avant que le sceau ne livre son mot de trois lettres...
		
		Gravée au sol, une inscription hante encore la sortie :
		\upshape
		\begin{center}
			« \textbf{La clé que tu cherches porte un nom de pierre.} »
		\end{center}
		\upshape
		
		\textbf{Dans cet épisode :} le Sceau d'émeraude va livrer son mot, trois
		gardiens menteurs veillent sur le Sanctuaire de jade, et le coffre
		n'attend plus que son code...
	\end{precedemment}
	
	\begin{conseil}
		\textbf{Acte I : rituel de réveil (5 min).} La crypte s'est rouverte.
		Sans relire l'épisode 1, recopie dans la marge, depuis ton carnet : la
		valeur de $5^{-1}$ modulo 7, ta matrice $K_1^{-1}$, le sceau $K_2$ et le
		vrai parchemin $K_2^{-1}$. Ce sont tes armes pour aujourd'hui.
	\end{conseil}
	
	\subsection*{\textcolor{emeraude}{Acte II : Le mot de l'émeraude (épreuve 2, fin)} \normalsize\textcolor{emeraude}{(10 min)} : 3 XP}
	
	\begin{enumerate}[label=\textbf{B.\arabic*.},leftmargin=10mm,start=3]
		\item \textbf{Le deuxième fragment.} Sur le Sceau d'émeraude sont gravés
		trois nombres premiers :
		\[
		\begin{pmatrix}53\\79\\59\end{pmatrix}.
		\]
		Appliquer le rituel R1 (module 13), puis décrypter avec la vraie clé
		$K_2^{-1}$ : chaque nombre obtenu code une lettre selon l'alphabet
		réduit ($\text{A}=0$, $\dots$, $\text{M}=12$; rituel R3).
		\textbf{Décrypter} le deuxième fragment.\Epx{3}
	\end{enumerate}
	
	% =====================================================================
	\subsection*{\textcolor{jade}{Acte III : Le Sanctuaire de jade (épreuve 3)} \normalsize\textcolor{jade}{(niveau 3 : matrices $4\times 4$, rituel modulo 13, 25 min)} : 5 XP}
	
	Au fond de la crypte, la Porte de jade ne répond qu'à son inverse exact.
	Gravée sur la porte, la matrice
	\[
	P=\begin{pmatrix}
		1&1&0&0\\
		0&1&1&0\\
		0&0&1&1\\
		2&0&0&1
	\end{pmatrix}.
	\]
	Trois \emph{gardiens} se présentent à toi, chacun prétendant être
	l'inverse de $P$ :
	\[
	G_1=\begin{pmatrix}
		12&1&12&1\\
		2&12&1&12\\
		11&2&12&1\\
		2&11&2&12
	\end{pmatrix}\qquad
	G_2=\begin{pmatrix}
		12&1&12&1\\
		12&12&1&12\\
		11&2&12&1\\
		2&11&2&12
	\end{pmatrix}\qquad
	G_3=\begin{pmatrix}
		12&1&12&1\\
		2&12&1&12\\
		2&2&12&1\\
		2&11&2&12
	\end{pmatrix}.
	\]
	
	\begin{enumerate}[label=\textbf{C.\arabic*.},leftmargin=10mm]
		\item \textbf{L'astuce du voleur.} En \textbf{comparant} les trois
		gardiens, le voleur constate qu'ils ne diffèrent \emph{que par leur
			première colonne}. Pas besoin d'un produit $4\times 4$ complet pour
		les départager : il suffit de vérifier que la première colonne du
		produit $P G$ est la première colonne de $I_4$, c'est-à-dire
		$\begin{pmatrix}1\\0\\0\\0\end{pmatrix}$.
		(Le voleur a une preuve que ce test ne peut rien lui cacher : elle
		attend les héros honnêtes dans le Carnet du Gardien.)
		\begin{enumerate}[label=\alph*)]
			\item Calculer $P$ multiplié par la première colonne de $G_1$,
			puis de $G_2$ et de $G_3$, modulo 13.\Epx{1,5}
			\item \textbf{Comparer} : quel gardien passe l'épreuve et peut
			prétendre être $P^{-1}$ ?\Epx{0,5}
		\end{enumerate}
		\item \textbf{L'honnêteté du héros.} Vérifier que
		$P G_1 \equiv I_4 \pmod{13}$ (détailler la première colonne, puis
		décrire la méthode pour les suivantes).\Epx{1}
		\item \textbf{Le dernier fragment.} Sur la porte de jade, quatre nombres
		premiers sont gravés :
		\[
		\begin{pmatrix}61\\3\\7\\113\end{pmatrix}.
		\]
		Appliquer le rituel R1 (module 13), puis décrypter avec le vrai gardien
		$G_1$ : chaque nombre obtenu code une lettre selon l'alphabet réduit
		($\text{A}=0$, $\dots$, $\text{M}=12$). \textbf{Décrypter} le dernier
		fragment.\Epx{2}
	\end{enumerate}
	
	% =====================================================================
	\subsection*{\textcolor{or}{Générique de fin : le coffre de l'Ordre (épilogue)} \normalsize\textcolor{or}{(10 min)} : 2 XP}
	
	\begin{missive}{La dernière énigme}
		\begin{enumerate}[label=\textbf{E.\arabic*.},leftmargin=10mm]
			\item Assembler les trois fragments décryptés. \textbf{Quel est le nom du
				trésor de l'Ordre ?} (Restituer les accents perdus par la
			Crypte.)\Epx{1}
			\item Le coffre ne s'ouvre qu'en entrant la \textbf{somme des nombres
				premiers de la Table de l'Ordre} correspondant aux neuf lettres du
			message. \textbf{Quel est le code du coffre ?}\Epx{1}
		\end{enumerate}
	\end{missive}
	
	\vspace{2mm}
	\begin{center}
		\itshape Le Grand Maître attend ton rapport complet, rédigé avec le soin
		d'un véritable cryptographe.\\[2mm]
		(Fin de l'épisode 2. Fin de la saison 1. La Crypte, elle, garde
		d'autres secrets...)
	\end{center}
	
	% =====================================================================
	\newpage
	\section*{Annexes de la Crypte}
	
	\subsection*{Annexe A : La Table de l'Ordre (nombres premiers)}
	\begin{center}
		\begin{tabular}{*{13}{c}}
			\toprule
			A & B & C & D & E & F & G & H & I & J & K & L & M\\
			\midrule
			2 & 3 & 5 & 7 & 11 & 13 & 17 & 19 & 23 & 29 & 31 & 37 & 41\\
			\midrule
			N & O & P & Q & R & S & T & U & V & W & X & Y & Z\\
			\midrule
			43 & 47 & 53 & 59 & 61 & 67 & 71 & 73 & 79 & 83 & 89 & 97 & 101\\
			\bottomrule
		\end{tabular}
	\end{center}
	
	\subsection*{Annexe B : Alphabet réduit (sceaux modulo 13)}
	Chaque nombre $x$ de $0$ à $12$ code une lettre :
	\begin{center}
		\begin{tabular}{l|*{13}{c}}
			& x=0 & 1 & 2 & 3 & 4 & 5 & 6 & 7 & 8 & 9 & 10 & 11 & 12\\
			\hline
			Lettre & A & B & C & D & E & F & G & H & I & J & K & L & M\\
		\end{tabular}
	\end{center}
	
	\subsection*{Annexe C : Alphabet complet par paires (sceau modulo 7)}
	Chaque lettre de rang $\ell$ ($\text{A}=0$, $\dots$, $\text{Z}=25$) est portée
	par le couple $(a,b)$ tel que $\ell = a + 7b$ :
	\begin{center}
		\begin{tabular}{*{13}{c}}
			\toprule
			A & B & C & D & E & F & G & H & I & J & K & L & M\\
			(0,0) & (1,0) & (2,0) & (3,0) & (4,0) & (5,0) & (6,0) & (0,1) & (1,1) & (2,1) & (3,1) & (4,1) & (5,1)\\
			\midrule
			N & O & P & Q & R & S & T & U & V & W & X & Y & Z\\
			(6,1) & (0,2) & (1,2) & (2,2) & (3,2) & (4,2) & (5,2) & (6,2) & (0,3) & (1,3) & (2,3) & (3,3) & (4,3)\\
			\bottomrule
		\end{tabular}
	\end{center}
	
	\subsection*{Annexe D : Mémo du cryptographe}
	\begin{itemize}
		\item \textbf{Produit :} si $M$ ($n$ lignes, $p$ colonnes) et $N$
		($p$ lignes, $q$ colonnes), le coefficient $(i,j)$ de $MN$ est
		$\sum_{k=1}^{p} m_{ik}\, n_{kj}$.
		\item \textbf{Inverse $2\times 2$ :} si $D = \begin{pmatrix}a&b\\c&d\end{pmatrix}$
		avec $\delta = ad-bc \neq 0$, alors
		$D^{-1} = \dfrac{1}{\delta}\begin{pmatrix}d&-b\\-c&a\end{pmatrix}$.
		Modulo $n$ : remplacer $\frac{1}{\delta}$ par l'inverse de $\delta$
		modulo $n$ (s'il existe).
		\item \textbf{Le théorème des voleurs de clés :}
		$(MN)^{-1} = N^{-1}M^{-1}$ : l'ordre s'inverse.
		\item \textbf{Inversibles modulo $n$ :} $x$ est inversible modulo $n$
		si et seulement si $x$ et $n$ sont premiers entre eux
		(Bézout). C'est pourquoi l'Ordre chérit les modules
		premiers : tout $x$ non nul y est inversible.
	\end{itemize}
	
	% =====================================================================
	\newpage
	\section*{\textcolor{argent}{Carnet du Gardien : éléments de correction}}
	
	\subsection*{Épisode 1 : La Pierre de Voûte (rappels, Sceau d'argent, parchemin du traître)}
	\subsection*{Épreuve 1 : Sceau d'argent}
	\textbf{A.1.} $K_1 = A + B - C
	= \begin{pmatrix}3+4-5 & 1+6-4\\ 2+0-1 & 5+4-5\end{pmatrix}
	= \begin{pmatrix}2&3\\1&4\end{pmatrix}$.\quad
	\textbf{A.2.} Comparaison coefficient par coefficient :
	seule $S_1 = \begin{pmatrix}2&3\\1&4\end{pmatrix}$ égale $K_1$
	($S_2$ est la transposée fausse, $S_3$ contient $-3$ au lieu de $3$).
	
	\textbf{A.3.} a) $\det K_1 = 2\times 4 - 3\times 1 = 5$.
	Comme $3\times 5 - 2\times 7 = 15 - 14 = 1$, on a
	$3\times 5 \equiv 1 \pmod 7$ : l'inverse de $5$ modulo 7 est $3$.
	b) $K_1^{-1} \equiv 3\begin{pmatrix}4&-3\\-1&2\end{pmatrix}
	= \begin{pmatrix}12&-9\\-3&6\end{pmatrix}
	\equiv \begin{pmatrix}5&5\\4&6\end{pmatrix} \pmod 7$
	(car $12\equiv 5$, $-9\equiv 5$, $-3\equiv 4$).
	Vérification :
	$K_1 K_1^{-1} \equiv \begin{pmatrix}2&3\\1&4\end{pmatrix}
	\begin{pmatrix}5&5\\4&6\end{pmatrix}
	= \begin{pmatrix}22&28\\21&29\end{pmatrix}
	\equiv \begin{pmatrix}1&0\\0&1\end{pmatrix} \pmod 7$. $\checkmark$
	
	\textbf{A.4.} a) Rituel R1 :
	$\begin{pmatrix}11&7\\29&7\end{pmatrix}
	\equiv \begin{pmatrix}4&0\\1&0\end{pmatrix} = Y \pmod 7$
	(les seuls premiers divisibles par 7 sont... 7 lui-même !).
	b) $X = K_1^{-1} Y \equiv \begin{pmatrix}5&5\\4&6\end{pmatrix}
	\begin{pmatrix}4&0\\1&0\end{pmatrix}
	= \begin{pmatrix}25&0\\22&0\end{pmatrix}
	\equiv \begin{pmatrix}4&0\\1&0\end{pmatrix} \pmod 7$.
	c) Première colonne : $(4,1)$ donne le rang $4 + 7\times 1 = 11$,
	soit la lettre L. Seconde colonne : $(0,0)$ donne le rang $0$, soit A.
	\textbf{Fragment 1 : « LA »}.
	
	\subsection*{Épreuve 2 : Sceau d'émeraude}
	\textbf{B.1.} $K_2 = MN = \begin{pmatrix}9&2&0\\4&16&3\\0&5&1\end{pmatrix}
	\equiv \begin{pmatrix}9&2&0\\4&3&3\\0&5&1\end{pmatrix} \pmod{13}$
	(car $16 \equiv 3$).
	
	\textbf{B.2.} a) $P_1 = M'N' = \begin{pmatrix}12&7&6\\3&1&10\\7&8&1\end{pmatrix}$
	et $P_2 = N'M' = \begin{pmatrix}1&11&6\\9&9&12\\7&7&6\end{pmatrix}$
	modulo 13. Ils sont \emph{différents} : le produit de matrices n'est pas
	commutatif.
	b) $K_2 P_1 \equiv \begin{pmatrix}10&0&9\\0&3&5\\9&0&12\end{pmatrix}
	\neq I_3$ : le parchemin $P_1 = M'N'$ est le \textbf{faux}.
	$K_2 P_2 \equiv I_3 \pmod{13}$ : le parchemin $P_2 = N'M'$ est le vrai,
	conformément au théorème $(MN)^{-1} = N^{-1}M^{-1}$.
	Donc $K_2^{-1} = \begin{pmatrix}1&11&6\\9&9&12\\7&7&6\end{pmatrix}$.
	
	\subsection*{Épisode 2 : Le Nom de Pierre (mot de l'émeraude, Sanctuaire de jade, coffre)}
	
	\textbf{B.3.} Rituel R1 : $53 \equiv 1 \pmod{13}$,
	$79 \equiv 1 \pmod{13}$, $59 \equiv 7 \pmod{13}$, donc
	$Y = \begin{pmatrix}1\\1\\7\end{pmatrix}$.
	Alors
	\[
	X = K_2^{-1}Y \equiv \begin{pmatrix}1&11&6\\9&9&12\\7&7&6\end{pmatrix}
	\begin{pmatrix}1\\1\\7\end{pmatrix}
	= \begin{pmatrix}54\\102\\56\end{pmatrix}
	\equiv \begin{pmatrix}2\\11\\4\end{pmatrix} \pmod{13},
	\]
	soit les lettres $2 \to$ C, $11 \to$ L, $4 \to$ E.
	\textbf{Fragment 2 : « CLE »}.
	
	\subsection*{Épreuve 3 : Sanctuaire de jade}
	\textbf{C.1.} a) Premières colonnes :
	$G_1 : (12,2,11,2)$, $G_2 : (12,12,11,2)$, $G_3 : (12,2,2,2)$.
	\[
	P\,(12,2,11,2)^{\!\top} \equiv (12+2,\ 2+11,\ 11+2,\ 24+2)^{\!\top}
	= (14,13,13,26)^{\!\top} \equiv (1,0,0,0)^{\!\top} \pmod{13}.
	\]
	\[
	P\,(12,12,11,2)^{\!\top} \equiv (24,23,13,26)^{\!\top}
	\equiv (11,10,0,0)^{\!\top} \neq (1,0,0,0)^{\!\top},
	\]
	\[
	P\,(12,2,2,2)^{\!\top} \equiv (14,4,4,26)^{\!\top}
	\equiv (1,4,4,0)^{\!\top} \neq (1,0,0,0)^{\!\top}.
	\]
	b) Seul $G_1$ passe l'épreuve du voleur : c'est lui qui prétend à juste
	titre être $P^{-1}$.
	
	\textbf{C.2.} La première colonne est détaillée en C.1 :
	$P\,(12,2,11,2)^{\!\top} \equiv (1,0,0,0)^{\!\top} \pmod{13}$.
	Colonne 2 de $G_1$ : $(1,12,2,11)^{\!\top}$, et
	$P\,(1,12,2,11)^{\!\top} = (13,14,13,13)^{\!\top}
	\equiv (0,1,0,0)^{\!\top}$. Colonne 3 : $(12,1,12,2)^{\!\top}$, et
	$P\,(12,1,12,2)^{\!\top} = (13,13,14,26)^{\!\top}
	\equiv (0,0,1,0)^{\!\top}$. Colonne 4 : $(1,12,1,12)^{\!\top}$, et
	$P\,(1,12,1,12)^{\!\top} = (13,13,13,14)^{\!\top}
	\equiv (0,0,0,1)^{\!\top}$. Toutes les colonnes sont exactes :
	$P G_1 \equiv I_4 \pmod{13}$.
	
	\emph{Pourquoi l'astuce du voleur fonctionne :} les trois gardiens ne
	diffèrent que par leur première colonne. Supposons qu'un candidat $G'$
	diffère de $G_1$ uniquement par sa première colonne, d'écart
	$\delta = (\delta_1,\delta_2,\delta_3,\delta_4)^{\!\top}$ : la première
	colonne de $P\,G'$ vaut alors
	\[
	e_1 + P\,\delta
	= \begin{pmatrix}1\\0\\0\\0\end{pmatrix}
	+ \begin{pmatrix}\delta_1+\delta_2\\ \delta_2+\delta_3\\
		\delta_3+\delta_4\\ 2\delta_1+\delta_4\end{pmatrix}.
	\]
	Le candidat passe le test du voleur si et seulement si $P\,\delta \equiv
	\vec{0} \pmod{13}$. Or $P\,\delta \equiv \vec{0}$ entraîne
	$\delta_2 \equiv -\delta_1$, $\delta_3 \equiv \delta_1$,
	$\delta_4 \equiv -\delta_1$, puis
	$2\delta_1 + \delta_4 \equiv 2\delta_1 - \delta_1 = \delta_1 \equiv 0$ :
	donc $\delta \equiv \vec{0}$ et le candidat était $G_1$ lui-même. Un faux
	gardien est donc toujours démasqué par la seule première colonne.
	
	\textbf{C.3.} Rituel R1 : $61 \equiv 9 \pmod{13}$,
	$3 \equiv 3 \pmod{13}$, $7 \equiv 7 \pmod{13}$,
	$113 \equiv 9 \pmod{13}$, donc $Y = (9,3,7,9)^{\!\top}$.
	Alors
	\[
	X = G_1 Y \equiv \begin{pmatrix}
		12&1&12&1\\2&12&1&12\\11&2&12&1\\2&11&2&12
	\end{pmatrix}
	\begin{pmatrix}9\\3\\7\\9\end{pmatrix}
	= \begin{pmatrix}204\\169\\198\\173\end{pmatrix}
	\equiv \begin{pmatrix}9\\0\\3\\4\end{pmatrix} \pmod{13},
	\]
	soit les lettres $9 \to$ J, $0 \to$ A, $3 \to$ D, $4 \to$ E.
	\textbf{Fragment 3 : « JADE »}.
	
	\subsection*{Épilogue}
	\textbf{E.1.} Les fragments s'assemblent : « LA » + « CLE » + « JADE ».
	Le trésor de l'Ordre est \textbf{la Clé de Jade}.
	\textbf{E.2.} Avec la Table de l'Ordre :
	L(37) + A(2) + C(5) + L(37) + E(11) + J(29) + A(2) + D(7) + E(11)
	$= 141$.
	\textbf{Le code du coffre est 141.}
	
	\vspace{3mm}
	\begin{conseil}
		\textbf{Prolongements possibles :} faire générer aux élèves leurs propres
		clés de chiffrement (matrice $2\times2$ de déterminant inversible modulo 7)
		et s'échanger des messages ; chercher ce qui se passe lorsque le
		déterminant n'est pas inversible modulo $n$ (message indéchiffrable,
		matrice non inversible dans $\Z/n\Z$) ; retrouver le message d'un camarade
		par cryptanalyse.
	\end{conseil}
	
\end{document}
